\section{Finite Approximations and the Stage Construction}

\subsection{Approximations and Limits}

The stage construction maintains lists of numbers already enumerated into
the two sets.  The finite-set approximations and limits are:

\begin{lstlisting}
def AstageF (s : Nat) : Finset Nat :=
  (stageState s).Alist.toFinset

def BstageF (s : Nat) : Finset Nat :=
  (stageState s).Blist.toFinset

def Aset : Set Nat := limitSet AstageF
def Bset : Set Nat := limitSet BstageF
\end{lstlisting}

The theorems \texttt{AstageF_mono} and \texttt{BstageF_mono} prove that
these approximations are monotone.  The generic approximation file proves
\texttt{StageMono.stabilizesBelow}: below every finite bound, a monotone
finite-stage approximation eventually agrees with its limit.  The
satisfaction proof uses this fact through
\texttt{run_halt_snapshot_of_limit}.

\subsection{Strategy State}

Each requirement has a record:

\begin{lstlisting}
structure ReqState where
  witness : Option Nat
  acted : Bool
  restraint : Nat
\end{lstlisting}

The construction state is:

\begin{lstlisting}
structure ConsState where
  Alist : List Nat
  Blist : List Nat
  reqs : List ReqState
  fresh : Nat
\end{lstlisting}

The \texttt{fresh} counter is intended to be strictly above enumerated
numbers and live witnesses and weakly above restraints.  This intention is
not left as a comment: it is part of the invariant \texttt{ConsInv} in
\texttt{Invariants.lean}.

\subsection{Attention and Action}

Even priority \(i=2e\) attacks \(A=\Phi_e^B\) and therefore queries the
current \(B\)-list.  Odd priority \(i=2e+1\) attacks \(B=\Phi_e^A\) and
queries the current \(A\)-list.  The definitions
\texttt{oracleOf} and \texttt{targetOf} implement this parity split.

The convergence check is:

\begin{lstlisting}
def convCheck (st : ConsState) (s i x : Nat) : Option Nat :=
  let v := nrun s (snapshotOf (st.oracleOf i) s) (i / 2) x
  if 2 <= v /\ (v - 2).unpair.1 = 0
  then some (v - 2).unpair.2 else none
\end{lstlisting}

It runs the numbered program \(i/2\) with fuel \(s\) against a length-\(s\)
snapshot of the current oracle side.  If the encoded run halts with output
\(0\), it returns the recorded use.

The attention predicate is:

\begin{lstlisting}
def requiresAttention (st : ConsState) (s i : Nat) : Bool :=
  match st.reqs[i]? with
  | none => false
  | some r =>
    match r.witness with
    | none => true
    | some x => !r.acted && (st.convCheck s i x).isSome
\end{lstlisting}

A requirement requires attention if its record exists and either no witness
has been appointed, or it has an unacted witness with visible output \(0\).
The deterministic stage step appends the newborn blank requirement record,
finds the least requirement requiring attention, and calls
\texttt{stepAt}.  Appointment stores the fresh counter as witness; action
enumerates the witness into the target side, records the use as restraint,
and replaces every lower-priority record by \texttt{ReqState.init}.

\subsection{Stage Dynamics and Invariants}

\texttt{StageDynamics.lean} defines \texttt{preState} and
\texttt{attended}, proves that \texttt{attended} is the least requirement
requiring attention, and packages the transition in \texttt{step_cases}.
It also proves \texttt{Alist_succ} and \texttt{Blist_succ}, which say that
each side either stays unchanged or receives exactly the acting witness.

\texttt{Invariants.lean} introduces the ten-field invariant
\texttt{ConsInv}.  Its fields include freshness of enumerated numbers,
freshness of witnesses, boundedness of restraints, coherence of blank
records, distinctness of live witnesses, witness-respecting higher
restraints, unacted witnesses outside both lists, acted witnesses in their
target list, and preservation of the acted computation.  The theorem
\texttt{consInv_stage} proves that every reachable \texttt{stageState s}
satisfies the invariant by induction over the three transition cases:
record birth, appointment, and action.

\begin{figure}
\centering
\begin{tikzpicture}[
  node distance=7mm,
  box/.style={draw, rounded corners, align=center, minimum width=35mm, minimum height=8mm, fill=blue!4},
  mbox/.style={draw, rounded corners, align=center, minimum width=35mm, minimum height=8mm, fill=orange!12},
  arr/.style={-Latex, thick}
]
\node[box] (oracle) {OracleCode\\FiniteEval};
\node[box, right=of oracle] (use) {Use\\InfiniteEval};
\node[box, below=of oracle] (approx) {Approximation};
\node[box, below=of use] (construct) {Construction\\StageDynamics};
\node[box, below=of construct] (inv) {Invariants\\FiniteInjury};
\node[box, below=of inv] (req) {Requirements};
\node[box, right=15mm of inv] (ce) {CE};
\node[box, below=of req] (main) {Main theorem};
\node[mbox, right=15mm of use] (mathlib) {Mathlib\\Primrec/Partrec};
\node[box, below=of mathlib] (bridge) {MathlibBridge\\RunPrimrec};

\draw[arr] (oracle) -- (use);
\draw[arr] (oracle) -- (approx);
\draw[arr] (use) -- (construct);
\draw[arr] (approx) -- (construct);
\draw[arr] (construct) -- (inv);
\draw[arr] (inv) -- (req);
\draw[arr] (req) -- (main);
\draw[arr] (construct) -- (ce);
\draw[arr] (bridge) -- (ce);
\draw[arr] (ce) |- (main);
\draw[arr] (mathlib) -- (bridge);
\end{tikzpicture}


\caption{Repository architecture.  Mathlib computability infrastructure
enters through the bridge and c.e. computability layer; the oracle
semantics and priority proof are project-local.}
\end{figure}

